\documentclass[11pt]{article}
\usepackage[margin=1in]{geometry}
\usepackage{amsmath,amssymb,amsthm,mathtools}
\usepackage{graphicx}
\usepackage{microtype}
\usepackage[hidelinks]{hyperref}
\usepackage{xcolor}

\newtheorem{theorem}{Theorem}
\newtheorem{lemma}{Lemma}
\newtheorem{corollary}{Corollary}
\theoremstyle{remark}
\newtheorem{remark}{Remark}

\newcommand{\C}{\mathbb C}
\newcommand{\T}{\mathbb T}
\newcommand{\kO}{\kappa_{O}^{\mathbb C}}
\newcommand{\kL}{\kappa_{L}^{\mathbb C}}

\title{\textbf{Roots-of-unity lower bounds for finite-dimensional\\
complex Littlewood constants}\\
\large A partial result for arXiv:1612.08397}
\author{Automated mathematics run \texttt{fa\_banach\_001}, lane 13}
\date{9 August 2026}

\begin{document}
\maketitle

\begin{center}
\fbox{\parbox{0.92\textwidth}{
\textbf{Status.} Candidate partial result; current verdict: \emph{likely
valid}. Human review is recommended. For every $n\ge2$ this packet gives an
explicit lower bound for both finite-dimensional complex constants in the
source problem and proves that their common value $1$ at $n=2$ cannot persist
for any $n\ge3$. It does not determine the exact constants beyond dimension
two.
}}
\end{center}

\section{Source question}

The source is W. Cavalcante, D. Pellegrino, and E. Teixeira,
\emph{On the geometry of multilinear forms}, arXiv:1612.08397, Section 5.2,
page 17 \cite{CPT}. For a bilinear form $T:\ell_\infty^n(\C)\times
\ell_\infty^n(\C)\to\C$, write $a_{ij}=T(e_i,e_j)$ and
\[
 \|T\|=\max_{|x_j|\le1,\ |y_i|\le1}
       \left|\sum_{i,j=1}^n a_{ij}y_i x_j\right|
 =\max_{|x_j|\le1}\sum_{i=1}^n\left|\sum_{j=1}^n a_{ij}x_j\right|.
\]
The optimal constants are defined by
\begin{align*}
 \left(\sum_{i=1}^n\left(\sum_{j=1}^n|a_{ij}|\right)^2\right)^{1/2}
   &\le \kO(n)\|T\|,\\
 \sum_{i=1}^n\left(\sum_{j=1}^n|a_{ij}|^2\right)^{1/2}
   &\le \kL(n)\|T\|.
\end{align*}
The source identifies their exact finite-dimensional values as an open
problem and proves in Theorem 25 that
\[
                    \kO(2)=\kL(2)=1.
\]

\begin{figure}[ht]
\centering
\includegraphics[width=\textwidth]{figures/open_problem_crop.png}
\caption{Definitions, open-problem statement, and the dimension-two theorem
on page 17 of the source paper.}
\end{figure}

\section{Proof intuition}

The most economical way to force a large mixed numerator is to give every
row two unimodular entries. Arrange their relative phases uniformly around
the circle: row $k$ is $(1,\zeta^k)$ for an $n$th root of unity $\zeta$.
The mixed numerator is then $n\sqrt2$, while the bilinear norm becomes the
maximum of a single equally spaced absolute-cosine sum.

That trigonometric maximum is exactly computable. On half a fundamental
period the signs of the cosines change only once. Removing the absolute
values therefore turns the sum into a signed geometric series. Its value is
real when $n$ is odd and has a simple phase shift when $n$ is even; both
cases have amplitude $\csc(\pi/(2n))$. Transposing the matrix preserves the
bilinear norm and converts the same example into one with equally large
Orlicz numerator.

\section{Result}

\begin{theorem}[Explicit roots-of-unity lower bounds]
\label{thm:main}
For every integer $n\ge2$,
\[
 \boxed{\quad
 \kO(n)\ \ge\ \frac{n}{\sqrt2}\sin\!\left(\frac{\pi}{2n}\right),
 \qquad
 \kL(n)\ \ge\ \frac{n}{\sqrt2}\sin\!\left(\frac{\pi}{2n}\right).
 \quad}
\]
Consequently, for every $n\ge3$ both constants are strictly greater than
$1$. In particular,
\[
               \kO(3),\ \kL(3)\ \ge\ \frac{3\sqrt2}{4}.
\]
The displayed lower bound increases with $n$ and converges to
$\pi/(2\sqrt2)$.
\end{theorem}

The proof uses the following exact elementary identity.

\begin{lemma}[Equally spaced absolute-cosine sum]
\label{lem:cosine}
For every integer $n\ge2$,
\[
 \max_{u\in\mathbb R}\sum_{k=0}^{n-1}
       \left|\cos\!\left(u+\frac{\pi k}{n}\right)\right|
     =\csc\!\left(\frac{\pi}{2n}\right).
\]
If $n$ is odd, a maximum occurs at $u=0$; if $n$ is even, one occurs at
$u=\pi/(2n)$.
\end{lemma}

\begin{proof}
Put $\alpha=\pi/n$ and
\[
                 S_n(u)=\sum_{k=0}^{n-1}|\cos(u+k\alpha)|.
\]
Reindexing the sum shows that $S_n$ is $\alpha$-periodic and even. It is
therefore enough to take $0\le u\le\alpha/2$.

First let $n=2m+1$. On this interval the cosines have positive sign for
$0\le k\le m$ and negative sign for $m+1\le k\le2m$ (with a harmless zero
at an endpoint). The geometric-sum identity gives
\[
 \sum_{k=0}^{m}e^{ik\alpha}-\sum_{k=m+1}^{2m}e^{ik\alpha}
       =\csc(\alpha/2).
\]
Taking the real part after multiplication by $e^{iu}$ yields
\[
                  S_n(u)=\csc(\alpha/2)\cos u.
\]
Its maximum on the chosen interval occurs at $u=0$.

Now let $n=2m$. The signs are positive for $0\le k\le m-1$ and negative for
$m\le k\le2m-1$. This time
\[
 \sum_{k=0}^{m-1}e^{ik\alpha}-\sum_{k=m}^{2m-1}e^{ik\alpha}
       =\cot(\alpha/2)-i.
\]
Hence
\begin{align*}
 S_n(u)
   &=\cot(\alpha/2)\cos u+\sin u\\
   &=\csc(\alpha/2)\cos(u-\alpha/2),
\end{align*}
whose maximum occurs at $u=\alpha/2$. Both cases prove the identity.
\end{proof}

\begin{proof}[Proof of Theorem~\ref{thm:main}]
Let $\zeta=e^{2\pi i/n}$ and let $A_n=(a_{kj})$ be the $n\times n$ matrix
whose only nonzero columns are
\[
             a_{k1}=1,\qquad a_{k2}=\zeta^k,
             \qquad 0\le k\le n-1.
\]
The second formula for the bilinear norm above gives
\[
 \|A_n\|=\max_{|x_1|,|x_2|\le1}
              \sum_{k=0}^{n-1}|x_1+\zeta^k x_2|.
\]
The objective is separately convex in $x_1$ and $x_2$, with each disk the
convex hull of its unit circle. Thus a maximum can be chosen with
$|x_1|=|x_2|=1$. After removing a common phase and writing
$x_2/x_1=e^{it}$, we obtain
\begin{align*}
 \|A_n\|
   &=\max_{t\in\mathbb R}\sum_{k=0}^{n-1}
          |1+e^{i(t+2\pi k/n)}|\\
   &=2\max_{u\in\mathbb R}\sum_{k=0}^{n-1}
          \left|\cos\!\left(u+\frac{\pi k}{n}\right)\right|\\
   &=2\csc\!\left(\frac{\pi}{2n}\right),
\end{align*}
where the last equality is Lemma~\ref{lem:cosine}.

Every row of $A_n$ has Euclidean norm $\sqrt2$. Therefore its mixed
Littlewood numerator is $n\sqrt2$, and
\[
 \kL(n)\ge\frac{n\sqrt2}{2\csc(\pi/(2n))}
          =\frac{n}{\sqrt2}\sin\!\left(\frac{\pi}{2n}\right).
\]

For the Orlicz inequality, use $A_n^{\mathsf T}$. Transposition preserves the
bilinear norm. The first two rows of $A_n^{\mathsf T}$ have $\ell_1$ norm
$n$, and all remaining rows vanish, so its Orlicz numerator is
$\sqrt{n^2+n^2}=n\sqrt2$. This proves the second lower bound.

At $n=2$ the bound is $1$. To see strict growth thereafter, set
$f(x)=x\sin(\pi/(2x))$ and $t=\pi/(2x)$. Then
\[
 f'(x)=\sin t-t\cos t>0,
\]
because $\sin t-t\cos t$ vanishes at $t=0$ and has derivative
$t\sin t>0$. Finally,
\[
 \lim_{n\to\infty}\frac{n}{\sqrt2}\sin\!\left(\frac{\pi}{2n}\right)
 =\frac{\pi}{2\sqrt2}.
\]
\end{proof}

\section{Verification, limitations, and novelty}

The script \texttt{code/verify\_roots\_of\_unity\_bound.py} samples a
fundamental half-period of $S_n$, checks the exact norm formula, and compares
the resulting matrix ratio with the theorem for $2\le n\le12$. This is only
numerical corroboration; the proof above is independent of it.

The theorem does not determine either exact finite-dimensional constant for
$n\ge3$. The first unresolved candidate equality is
\[
              \kO(3)=\kL(3)=\frac{3\sqrt2}{4}.
\]
A bounded attempt to prove the necessary upper bound found no valid extension
of the source's dimension-two classification. The exact missing inequality
is documented in the accompanying dimension-three frontier attempt note.

A later paper of Cavalcante, N\'u\~nez-Alarc\'on, Pellegrino, and Rueda proves
the dimension-free complex mixed constant $2/\sqrt\pi$ \cite{CNPR}. Thus the
present asymptotic lower bound $\pi/(2\sqrt2)$ lies below the known
dimension-free optimum, as it must. That later result does not provide the
finite-$n$ formula requested by the source.

The bounded novelty audit on 9 August 2026 searched the arXiv full text and
citation neighborhood of arXiv:1612.08397, exact variants of the displayed
formula, and combinations of ``roots of unity,'' ``complex Littlewood,'' and
``finite-dimensional constants.'' It also checked arXiv:1912.10313. No match
for this construction or all-$n$ lower bound was located. Novelty is therefore
provisional, not certified.

\paragraph{Human-review recommendation.}
The result is short enough for direct verification. Review should focus on
the two signed geometric-sum identities in Lemma~\ref{lem:cosine} and on the
separate-convexity reduction from the bidisk to the torus.

\begin{thebibliography}{9}

\bibitem{CPT}
W. Cavalcante, D. Pellegrino, and E. Teixeira,
\emph{On the geometry of multilinear forms},
arXiv:1612.08397 (2016), Section 5.2.

\bibitem{CNPR}
W. Cavalcante, D. N\'u\~nez-Alarc\'on, D. Pellegrino, and P. Rueda,
\emph{Optimal constants of classical inequalities for complex scalars},
arXiv:1912.10313 (2019).

\end{thebibliography}

\end{document}
